
\section{Discussion}

As show analytically in Section \ref{sec:adv_loss}, in the optimal case adversarial training can perform as well as our derived method; it is also intuitively simple and allows for more nuanced tuning. However, it introduces an extra layer of complexity (indeed, a second optimization problem) into our system. In this particular case of invariant representation, our results lead us to believe that adversarial training is unnecessary.

This does not mean that adversarial training for invariant representations is strictly worse in practice. There are certainly cases where training an adversary may be easier or less restrictive than other methods, and due to its shared literature with Generative Adversarial Networks \cite{goodfellow2014generative}, there may be training heuristics or other techniques that can improve performance.

On the otherhand, we believe that our derivations here shed light on why these methods might fail. We believe specific failure modes of adversarial training can be attributed to Eq. \ref{eq:adv_inf}, where the adversary fails to achieve the infimum. Bad approximations (i.e. weak or poorly trained adversaries) may provide bad gradient information to the system, leading to poor performance of the encoder against a post-hoc adversary.

%\subsection{Note on the experimental results of Xie et al.}

Our experimental results do not match those reported in Xie et al. While in general their method has comparable performance for predictive accuracy, we do not find that their adversarial error is low; instead, we find that the encoder/adversary pair becomes stuck in local minima.
%We believe this is reflected by the t-SNE plot.
We also find that the adversary trained alongside the encoder performs badly against the encoder (i.e. the adversary cannot predict $c$ well), but a post-hoc trained adversary performs very well, easily predicting $c$ (as demonstrated by our experiments).

It may be that we have inadvertently built a stronger adversary. We have attempted to follow the author's experimental design as closely as possible, using the same architecture and the same adversary (using the gradient-flip trick and 3-layer feed forward networks).
%We believe the adversarial case has the potential for improvement with careful tuning, and perhaps the removal of the gradient-flip trick.
%\BlueComment{, particularly by modifying their architecture and detaching the training of the adversary from the encoder (i.e. removing the gradient-flip trick), which provides greater degrees of freedom to tune the adversary}
With the details provided we could not replicate their reported adversarial error for their method, nor for the VFAE method. However, we are able to reproduce the adversarial error reported in Louizos et al., which uses logisic regression. In general for stronger adversaries the adversarial loss will increase, but the relative rankings should remain roughly the same.

%\subsection{Correlation of $c$ and $y$}




% \subsection{Invariance and Equivariance vs. Independence}

% The concept of invariant or equivariant features are actually slightly different formally than the independence condition $z \perp c$. In the invariance frame, $c$ is viewed as a group or semi-group of transformations (e.g. rotations). Encoding $z$ is said to be invariant to $c$ if $p(z|x) = p(z|c(x))$ for all $c$ and $x$. We can draw a rough equivalence to our context by defining $\tilde{x}$ to be the quotient representation $x\backslash c$. In the optimal case in our context, $I(z,c) = 0$, which implies that $p(z|c)=p(z)$. The following also then holds:
% \begin{align}
% \int p(\tilde{x}) p(z|\tilde{x},c) = p(z,c) = p(z) = \int p(\tilde{x}) p(z|\tilde{x})
% \end{align}
% Thus, $p(z|\tilde{x},c) = p(z|\tilde{x})$ almost everywhere, which is the invariance condition. Computationally, however, such formalism is unnecessary, as the resulting condition is the same ($z$ uninformed of $c$).


\section{Conclusion}

We have derived a variational upper bound for the mutual information between latent representations and covariate factors. Provided a dataset with labeled covariates, we can train both supervised and unsupervised learning methods that are invariant to these factors without the use of adversarial training. After training our method can be used in production without requiring covariate labels. Finally, our approach also enables manipulation of specified factors when generating realistic data. Our direct, information-theoretic optimization approach avoids the pitfalls inherent in adversarial learning for invariant representation and produces results that match or exceed capabilities of these state-of-the-art methods.

\subsection*{Acknowledgements}

This work was supported by DARPA grants W911NF-16-1-0575 and FA8750-17-C-0106, as well as the NSF Graduate Research Fellowship Program Grant Number DGE-1418060. We would like to thank the conference organizers, area chairs, and especially the anonymous reviewers for their work and helpful input. We also would like to thank Ayush Jaiswal for several insightful conversations, Umang Gupta for finding and correcting an error in our Gaussian approximation, and Ishaan Gulrajani for finding and correcting a bug in our evaluation code.


